\clearpage\begingroup\SourceLeaf{345}\ArithmeticHead{14}{L}{STÄCKEL, GAUSS ALS GEOMETER.}
\renewcommand{\thefootnote}{\arabic{footnote})}
\noindent im Leben müsste eine andere sein, wenn dergleichen Widerwärtigkeiten nicht
öfter eintreffen sollten«.
\par
Aus den vorstehenden Äusserungen klingt heraus, dass es nicht nur Mangel
an Musse war, der den Fortgang der theoretischen Forschungen hemmte, sondern
dass in der Gemütsverfassung von \textsc{Gauss} Hinderungen lagen. »Es ist
wahr«, schreibt er am 20. April 1848 an seinen Jugendfreund \textsc{Bolyai} (Br. G.-\textsc{Bolyai},
S. 132), »mein Leben ist mit Vielem geschmückt gewesen, was die
Welt für beneidenswert hält. Aber glaube mir, lieber \textsc{Bolyai}, die \textls[100]{herben}
Seiten des Lebens, wenigstens des meinigen, die sich wie der rote Faden dadurch
ziehen und denen man im höheren Alter immer wehrloser gegenübersteht,
werden nicht zum hundertsten Teil aufgewogen von dem Erfreulichen.
Ich will gern zugeben, dass dieselben Schicksale, die zu tragen mir so schwer
geworden ist und noch ist, manchem andern viel leichter gewesen wären, aber
die Gemütsverfassung gehört zu unserm Ich, der Schöpfer unserer Existenz
hat sie uns mitgegeben, und wir vermögen wenig daran zu ändern«. Es ist
hier nicht der Ort, von dem Leid zu sprechen, das \textsc{Gauss} mehr als einmal
in seinem Hause betroffen hat. Es hat »die Heiterkeit des Geistes«, die er
zur wissenschaftlichen Arbeit nötig hatte, »nur zu sehr und zu vielfach getrübt«
(Brief an \textsc{Bessel} vom 28. Februar 1839, Br. G.-\textsc{Bessel}, S. 524).
\par
Die Empfindlichkeit für unangenehme Eindrücke, von der \textsc{Gauss} in dem
Brief an \textsc{Bessel} vom 14. März 1824 spricht, hat sicherlich dazu beigetragen,
dass er es vermied, in seinen Veröffentlichungen Gegenstände zu berühren,
die zu Streitigkeiten Anlass geben konnten. Wie behutsam geht er in seiner
Dissertation mit den imaginären Grössen um, und gar ihre geometrische Deutung,
die er nach seinem Zeugnis schon vor 1799 besass, hat er damals unterdrückt
und erst 1831 bekannt gemacht. Ebenso hat er seine antieuklidische
Geometrie nicht zur Veröffentlichung ausgearbeitet. »Vielleicht wird dies auch
bei meinen Lebzeiten nie geschehen, da ich das Geschrei der Böoter scheue,
wenn ich meine Ansicht \textls[100]{ganz} aussprechen wollte« (Brief an \textsc{Bessel} vom
27. Januar 1829, W. VIII, S. 200).
\par
Diese Scheu war verstärkt worden durch böse Erfahrungen, die \textsc{Gauss}
machen musste, als er 1816 in der Besprechung der Parallelentheorien von
\textsc{Schwab} und \textsc{Metternich} (W. IV, S. 364, VIII, S. 170) Andeutungen über die
Unbeweisbarkeit des elften \textsc{Euklid}ischen Axioms gewagt hatte: »Es ist mit
\clearpage\SourceLeaf{346}\ArithmeticHead{15}{R}{DIE GRUNDLAGEN DER GEOMETRIE.}\setcounter{footnote}{0}
\noindent Kot danach geworfen«, schreibt er am 15. Januar 1827 an \textsc{Schumacher} (Br.
G.-\textsc{Sch.} II, S, 94)\footnote{Von wem der bösartige Angriff ausgegangen ist, hat sich noch nicht ermitteln lassen.}.
Solche Angriffe hatte \textsc{Gauss} wohl im Auge, wenn er
am 25. August 1818 an \textsc{Gerling} schrieb: »Ich freue mich, dass Sie den Mut
haben, sich [in Ihrem Lehrbuch] so auszudrücken, als wenn Sie die Möglichkeit,
dass unsere Parallelentheorie, mithin unsere ganze Geometrie, falsch
wäre, anerkennten. Aber die Wespen, deren Nest Sie aufstören, werden
Ihnen um den Kopf fliegen« (W. VIII, S. 179).
\par
Dazu kam die geringe Meinung, die \textsc{Gauss} von der grossen Mehrzahl der
Mathematiker hatte. Bereits am 16. Dezember 1799 schreibt er an \textsc{Wolfgang
Bolyai}, der ihm einen Versuch, das Parallelenaxiom zu beweisen, übersandt
hatte: »Mach’ doch ja Deine Arbeit bald bekannt; gewiss wirst Du dafür den
Dank zwar nicht des grossen Publikums (worunter auch mancher gehört, der
für einen geschickten Mathematiker gehalten wird) einernten, denn ich überzeuge
mich immer mehr, dass die Zahl der wahren Geometer äusserst gering
ist, und die meisten die Schwierigkeiten bei solchen Arbeiten weder beurteilen
noch selbst einmal sie verstehen können — aber gewiss den Dank aller derer,
deren Urteil Dir allein wirklich schätzbar sein kann« (W. VIII, S. 159). Als
\textsc{Wolfgang Bolyai} dann im Jahre 1832 seinem Jugendfreunde die \textit{Scientia
spatii absolute vera} seines Sohnes \textsc{Johann} übersandt hatte, in der das Rätsel
der Parallelenfrage gelöst war, antwortete dieser am 6 März 1832: »Die
meisten Menschen haben gar nicht den rechten Sinn für das, worauf es dabei
ankommt, und ich habe nur wenige Menschen gefunden, die das, was ich
ihnen mitteilte, mit besonderem Interesse aufnahmen. Um das zu können,
muss man erst recht lebendig gefühlt haben, was eigentlich fehlt, und darüber
sind die meisten Menschen ganz unklar« (W. VIII, S. 221). Noch schärfer
äussert sich \textsc{Gauss} in einem Briefe an \textsc{Gerling} vom 25. Juni 1815: »Mir däucht,
es ist in mehr als einer Rücksicht wichtig, bei den Schülern den Sinn für
Rigor wach zu erhalten, da die meisten Menschen nur gar zu geneigt sind,
zu einer laxen Observanz überzugehen. Selbst unsere grössten Mathematiker
haben meistenteils in dieser Rücksicht etwas stumpfe Fühlhörner« (Brief im
\textsc{Gauss}-Archiv). Ein gut Teil Menschenverachtung aber steckt in dem Rat, den
\textsc{Gauss} am 29. September 1837 seinem jüngeren Freunde \textsc{Möbius} erteilt: »Man
\clearpage\SourceLeaf{347}\ArithmeticHead{16}{L}{STÄCKEL, GAUSS ALS GEOMETER.}\setcounter{footnote}{0}
\noindent muss immer bedenken, dass, wo die Leser, für welche man schreibt, keinen
Anstoss nehmen, es vielleicht gar nicht wohlgetan wäre, tiefer einzudringen,
als ihnen frommt« (W. XI\textsubscript{1}, S. 19).
\par
\textsc{Gauss} hat bei seinen Klagen über mangelndes Verständnis wohl auch an
die Briefe gedacht, die er im Jahre 1831 mit \textsc{Schumacher} gewechselt hatte,
als dieser glaubte, das Parallelenaxiom bewiesen zu haben (W. VIII, S. 210—219).
\textsc{Schumacher} liess sich von der Unzulänglichkeit seines Verfahrens nicht
überzeugen und sandte den ausführlichen Brief \textsc{Gauss}ens vom 12. Juli 1831
an \textsc{Bessel}. »Eine tolle Geschichte«, antwortete dieser am 1. Aug. 1831, »ist
doch die im \textsc{Gauss}schen (hier zurückfolgenden) Briefe vorkommende, dass die
Peripherien zweier Kreise von den Halbmessern \(r\) und \(r'\) nicht im Verhältnis
\(r:r'\) stehen sollen. Ich bezweifele dieses nicht, weil \textsc{Gauss} es sagt; allein
diese Ungleichheit ist mir so wenig anschaulich, dass ich mir, nach dem alten
\textsc{Kulenkamp}schen Ausdruck\footnote{\textsc{Andreas Gottlieb Kulenkamp} hiess der Inhaber des Handelshauses in Bremen, bei dem \textsc{Bessel} von 1799 bis 1806 tätig gewesen war.} kein Denkbild davon machen kann« (Abschrift
des Briefes im \textsc{Gauss}-Archiv).
\par
Die Zurückhaltung, die \textsc{Gauss} übte, brachte die Gefahr mit sich, dass
andere ihm zuvorkamen, und das ist auch wiederholt geschehen. Aber in
diesem Punkte war \textsc{Gauss} unempfindlich. Am 30. Januar 1812 schreibt er
an \textsc{Laplace}: »J’ai dans mes papiers beaucoup de choses dont peut-être je
pourrai perdre la priorité de la publication, mais soit, j’aime mieux, faire
\textls[100]{mûrir} les choses« (W. X\textsubscript{1}, S. 374), und als \textsc{Abel} seine \textit{Recherches} veröffentlicht
hatte, begnügt er sich damit festzustellen, dass der Norweger ihn in
Bezug auf etwa ein Drittel der Sachen der Mühe überhoben habe, sie auszuarbeiten,
»zumal da er alle Entwicklungen mit vieler Eleganz und Konzision
gemacht« habe (Brief an \textsc{Bessel} vom 30. März 1828, Br. G.-\textsc{Bessel},
S. 477). Bei \textsc{Johann Bolyais} \textit{Scientia spatii} fand er es sogar höchst erfreulich,
dass gerade der Sohn seines alten Freundes ihm auf eine so merkwürdige Art
zuvorgekommen sei (Brief vom 6. März 1832, W. VIII, S. 221).
\clearpage\endgroup
